% SEGMENT OLP-0539-S01
% Part: set-theory
% Chapter: z
% Section: separation

\documentclass[../../../include/open-logic-section]{subfiles}

\begin{document}

\olfileid{sth}{z}{sep}
\olsection{பிரிப்பு}

முறைசாராப் பண்புவழிக் கணமாக்கலுக்குப் பதிலாகப் பின்வரும் கொள்கையுடன்
தொடங்குகிறோம்:

\begin{axiom}[பிரிப்பு அடிகோள் திட்டம்] ஒவ்வொரு வாய்பாடு $\phi(x)$ க்கும்,
பின்வருவது ஓர் அடிகோள்: எந்த $A$ க்கும்,
$\Setabs{x \in A}{\phi(x)}$ என்ற கணம் உள்ளது.
\end{axiom}

இது ஒரே ஓர் அடிகோள் அல்ல என்பதைக் கவனிக்கவும். இது அடிகோள்களின் ஓர்
\emph{திட்டம்}. \emph{முடிவிலாப் பல} பிரிப்பு அடிகோள்கள் உள்ளன;
ஒவ்வொரு வாய்பாடு $\phi(x)$ க்கும் ஒன்று உள்ளது. இத்திட்டத்தைப் பின்வரும்
வடிவிலும் அதேபோல் எழுதலாம் (வழக்கமாக அவ்வாறே எழுதப்படுகிறது):

\begin{defish}
``$S$'' ஐக் கொண்டிராத எந்த வாய்பாடு $\phi(x)$ க்கும், பின்வருவது ஓர்
அடிகோள்:
\[
	\forall A \exists S \forall x(x \in S \liff (\phi(x) \land x \in A)).
\]
\end{defish}

\olref[sth][][]{part} இன் தொடக்கத்தில் குறிப்பிட்ட மரபுக்கு இணங்க,
பிரிப்பு அடிகோள்களிலுள்ள வாய்பாடுகள்~$\phi$ அளவுருக்களைக்
கொண்டிருக்கலாம்.\footnote{இதன் பொருள் என்ன என்பதற்கான விளக்கத்திற்கு,
\olref[sfr][infinite][induction]{natinductionschema} ஐ அடுத்து உடனே வரும்
விவாதத்தைப் பார்க்கவும்.}

நாம் விளக்கிவரும் கணங்களின் திரள்வுறும் படிநிலைமுறைக் கருத்தாக்கம்
பிரிப்பை உடனடியாக நியாயப்படுத்துகிறது. ஏன் என்பதைப் பார்க்க, $A$ ஒரு
கணம் என்க. எனவே ஏதோ ஒரு படிநிலை~$S$ ஐ அடைவதற்குள் $A$ உருவாகிறது
(\stageshier{} இன்படி). $A$ படிநிலை~$S$ இல் உருவானதால், $A$ இன் அனைத்து
உறுப்புகளும் படிநிலை~$S$ க்கு முன்பு உருவாகின (\stagesacc{} இன்படி).
இப்போது குறிப்பாக, $A$ இன் உறுப்புகளாகவும் $\phi$ ஐ நிறைவேற்றுபவையாகவும்
உள்ள எல்லாக் கணங்களையும் கருதுக; இக்கணங்கள் அனைத்தும் படிநிலை~$S$ க்கு
முன்பு உருவாயின என்பது தெளிவு. ஆகவே அவையும் படிநிலை~$S$ இல்
$\Setabs{x \in A}{\phi(x)}$ என்ற கணமாக உருவாக்கப்படுகின்றன
(\stagesacc{} இன்படி).

முறைசாராப் பண்புவழிக் கணமாக்கலைப் போலன்றி, இது ரஸ்ஸலின் முரண்பாட்டைத்
தவிர்க்கிறது. ஏனெனில் $\Setabs{x}{x \notin x}$ என்ற கணம் உள்ளது என நாம்
வெறுமனே உறுதிப்படுத்த முடியாது. மாறாக, ஏதோ ஒரு கணம்~$A$ \emph{தரப்பட்டால்},
$R_A = \Setabs{x \in A}{x \notin x}$ என்ற கணம் உள்ளது என
உறுதிப்படுத்தலாம். ஆனால் இதனால் நிறுவப்படுவது $R_A \notin R_A$ மற்றும்
$R_A \notin A$ என்பவை மட்டுமே; அவற்றில் எதுவும் கவலைக்குரியதல்ல.
% SOURCE CORRECTION TA-STH-006: repair the source's subject--verb error “this avoid”.

எனினும், பிரிப்பிற்கு உடனடியானதும் வியப்பூட்டுவதுமான ஒரு விளைவு உண்டு:

\begin{thm}\ollabel{thm:NoUniversalSet}
\emph{அனைத்துமைக்} கணம் எதுவும் இல்லை; அதாவது,
$\Setabs{x}{x = x}$ என்பது இல்லை.
\end{thm}

\begin{proof}
முரண்பாட்டுவழி நிறுவலுக்காக, $V$ ஓர் அனைத்துமைக் கணம் எனக் கொள்க.
அப்போது பிரிப்பின்படி,
$R = \Setabs{x \in V}{x \notin x} = \Setabs{x}{x \notin x}$ உள்ளது;
இது ரஸ்ஸலின் முரண்பாட்டிற்கு முரணாகும்.
\end{proof}

அனைத்துமைக் கணம் இல்லாமை---உண்மையில், கணங்களின் படிநிலை வரிசை
முடிவில்லாமல் திறந்திருப்பது---திரள்வுறும் படிநிலைமுறைக் கருத்தாக்கத்தின்
மிக அடிப்படையான எண்ணங்களில் ஒன்று. ஆகவே, உள்ளுணர்வின்படி, வேறொரு வழியிலும்
இந்த முடிவை அடையலாம் என்பதைக் காண்பது பயனுள்ளது. அனைத்துமைக் கணம் தன்னுடைய
!!a{element} ஆக இருக்க வேண்டும். ஆனால் நமது திரள்வுறும் படிநிலைமுறைக்
கருத்தாக்கத்தின்படி, ஒவ்வொரு கணமும் அதன் !!{element}s அனைத்தும் தோன்றிய
படிநிலைக்கு உடனே அடுத்த முதல் படிநிலையில் (முதன்முறையாக) தோன்றுகிறது.
இதனால் \emph{எந்தக்} கணமும் தன்னுறுப்புடையதல்ல என்பது பின்வருகிறது.
ஏனெனில் தன்னுறுப்புடைய எந்தக் கணமும், அது முதன்முதலில் தோன்றிய
படிநிலைக்கு உடனே அடுத்த படிநிலையில்தான் முதலில் தோன்ற வேண்டும்; இது
முரணானது. (இந்த விளக்கத்தை இன்னும் துல்லியமாக்க, ஒரு கணத்தின் ``தரநிலை''
என்ற கருத்தை எவ்வாறு பயன்படுத்துவது என்பதை
\olref[sth][spine][rank]{defnsetrank} இல் காண்போம். ஆனால் அதற்கு முன்
மேலும் சில அடிகோள்கள் தேவைப்படும்.)

பிரிப்பினதும் உறுப்புசார் சமத்துவத்தினதும் மேலும் சில விளைவுகள் இதோ.

\begin{prop}\ollabel{prop:emptyexists}
ஏதாவது ஒரு கணம் இருந்தால், $\emptyset$ உள்ளது.
\end{prop}

\begin{proof}
$A$ ஒரு கணம் எனில், பிரிப்பின்படி
$\emptyset = \Setabs{x \in A}{x \neq x}$ உள்ளது.
\end{proof}

\begin{prop}
எந்தக் கணங்கள் $A$, $B$ க்கும் $A \setminus B$ உள்ளது.
\end{prop}

\begin{proof}
பிரிப்பின்படி $A \setminus B = \Setabs{x \in A}{x \notin B}$ உள்ளது.
\end{proof}

(கிட்டத்தட்ட) தன்னிச்சையான வெட்டுகளும் உள்ளன என்பது பின்வருகிறது:

\begin{prop}\ollabel{prop:intersectionsexist}
$A \neq \emptyset$ எனில்,
$\bigcap A = \Setabs{x}{(\forall y \in A)x \in y}$ உள்ளது.
\end{prop}

\begin{proof}
$A \neq \emptyset$ என்க; ஆகவே ஏதோ ஒரு $c \in A$ உள்ளது. அப்போது
$\bigcap A = \Setabs{x}{(\forall y \in A)x \in y} = \Setabs{x \in c}{(\forall y \in
A)x \in y}$; இது பிரிப்பின்படி உள்ளது.
\end{proof}

ஆனால் $A \neq \emptyset$ என்ற நிபந்தனையைக் கவனிக்கவும்; ஏனெனில்
$\bigcap \emptyset$ வெறுமைவழியில் அனைத்துமைக் கணமாக இருக்கும்; இது
\olref{thm:NoUniversalSet} க்கு முரணாகும்.

\end{document}
